\section{Enhanced NRB \& ORB (ENRB \& EORB)}

\subsection{Product Description}
\label{sec:enrb_eorb_product_description}
As is well known, SAR images are subject to various artifacts and noise sources that can hinder the implementation of many high-level applications. The Enhanced \acs{nrb} and \acs{orb} products are designed to retain the information content of their corresponding standard products, while mitigating these disturbances to improve overall image quality and usability.

These "enhanced" versions of the \acs{nrb} and \acs{orb} \acs{ard} products are generated by applying specific filtering procedures to the input \acs{slc} or \acs{cslc} data:

\begin{itemize}
    \item \textit{Ghost Removal:} Ghosts are spurious image artifacts caused by SAR ambiguities in the range or azimuth directions. These artifacts can significantly degrade application performance, particularly in tasks such as object detection and recognition. Ghost removal is especially beneficial for applications that utilize the \acs{orb} product.
    \item \textit{Sidelobe Suppression:} Sidelobes in SAR imagery may introduce two critical issues. First, \textit{false alarms} can occur when strong scatterers produce sidelobes that are mistakenly interpreted as genuine targets. Second, \textit{missed detections} may arise when the sidelobes of a strong target obscure the mainlobe of a nearby weaker target. Sidelobe suppression improves the reliability of detection, particularly in scenarios requiring either the \acs{orb} or \acs{nrb} products.
\end{itemize}

With the exception of the intensity cross-spectra, the Enhanced \acs{nrb} and \acs{orb} products closely match their standard counterparts (see \autoref{sec:nrb_prod_description} and \autoref{sec:orb_product}), but with additional processing steps aimed at reducing ghosting and sidelobe interference as much as possible.

The proposed intensity cross-spectra component follows the specification of Harmony’s Intensity Cross-Spectral Matrix (ICSM) product. For each polarization, the ICSM provides both intensity power spectra and cross-spectra, enhancing the frequency-domain analysis of the SAR backscatter.

\subsubsection{Input and Auxiliary Data}
The following input and auxiliary data are considered for product generation:

\begin{longtblr}[
    label = tab:in_aux_enh_products,
    entry = none,
    caption = {Enhanced \acs{nrb} and \acs{orb} input and auxiliary data},
    ]{
        width = \linewidth,
        colspec = {Q[l,20]Q[l,60]Q[c,20]},
        hlines,
        hline{1-Z} = {-}{0.05em},
    }
    \textbf{Name}   & \textbf{Description}   & \textbf{Mandatory (M) / Optional (O)} \\
    \acs{slc} or \acs{cslc} & L1 Single Look Complex image or co-registered L1 Single Look Complex image & M \\
    \acs{slc} annotation & L1 Single Look Complex image annotations or reference geometry annotations & M \\
    \acs{dem} & External DEM (e.g., Copernicus DEM at different resolutions, with 30 m as the default worldwide). & M for \acs{nrb} \\
    Precise Orbits & \acs{slc} precise orbits & M \\
    \acs{etad} & \acs{etad}-like products for atmospheric delay correction & O \\
    \acs{tec} & External \acs{tec} data for delay correction & O \\
    Geoid & External geoid & M for \acs{orb} \\
    \acs{era5} & ECMWF \acs{era5} wind data for angle of incidence compensation & M for \acs{orb} \\
\end{longtblr}

\subsubsection{Main Variables}
\label{sec:variables_enh}
The following tables list the main variables associated with the \acs{enrb} product.

\begin{longtblr}[
    label = tab:var_enrb_products,
    entry = none,
    caption = {\acs{enrb} main variables},
    ]{
        width = \linewidth,
        colspec = {Q[l,25]Q[l,52]Q[l,20]Q[l,13]},
        hlines,
        hline{1-Z} = {-}{0.05em},
    }
    \textbf{Name} & \textbf{Description} & \textbf{Dimensions} & \textbf{Data Type} \\
    gamma0 & Radiometric Terrain Correction gamma naught backscatter linear scaling image. Only for \acs{enrb} & polarization, acquisition\_date, x, y & float32 \\
    sigma0\_raw ($\sigma_{0,\mathrm{raw}}$) & Full resolution (single look) radar intensity, in $\sigma_0$ convention. Only for \acs{eorb}  & polarization, acquisition\_date, x, y & float32 \\
    sigma0\_nesz ($\sigma_{0,\mathrm{nesz}}$) & Spatially-averaged, LxL m$^2$ resolution noise equivalent sigma zero, in $\sigma_0$ convention. Only for \acs{eorb} & polarization, acquisition\_date, x, y & float32 \\
    sigma0\_denoise ($\sigma_{0,\mathrm{denoise}}$) & Spatially-averaged, LxL m$^2$ resolution denoized radar intensity, in $\sigma_0$ convention. Only for \acs{eorb} & polarization, acquisition\_date, x, y & float32 \\
    sigma0\_L ($\sigma_{0,L}$) & Spatially-averaged, LxL m$^2$ resolution denoized and outlier removed radar intensity, in $\sigma_0$ convention. Only for \acs{eorb} & polarization, acquisition\_date, x, y & float32 \\
    sigma0\_Lnorm ($\sigma_{0,L}/\sigma_{0,\mathrm{mod}}$) & Spatially-averaged, LxL m$^2$ resolution (mean) wind normalized backscatter measurements. Only for \acs{eorb} & polarization, acquisition\_date, x, y & float32 \\
    local\_inc\_angle & DEM-based local Incidence angle image [deg], only for \acs{enrb}. This variable is present only if the \acs{enrb} is generated from \acs{slc}s. If the \acs{enrb} is derived from \acs{cslc}s, this variable will be null, and the reference\_inc\_angle metadata will be present & x, y & float32 \\
    anf & Optional variable: area normalization factor or local contributing area, used to obtain the gamma0 coefficient; only for \acs{enrb}. This variable is present only if the \acs{enrb} is generated from \acs{slc}s. If the \acs{nrb} is derived from \acs{cslc}s, this variable will be null, and the reference\_anf metadata will be present & x, y & float32 \\
    mask & Full resolution (single look) bright target mask. Only for \acs{eorb} & polarization, acquisition\_date, x, y & bool \\
\end{longtblr}

The table below describes the dimensions/coordinates associated with the \acs{enrb} and \acs{eorb} main variables.

\begin{longtblr}[
    label = tab:var_enh_coordinates,
    entry = none,
    caption = {\acs{enrb} and \acs{eorb} main variables coordinates/dimensions},
    ]{
        width = \linewidth,
        colspec = {Q[l,20]Q[l,62]Q[l,18]},
        hlines,
        hline{1-Z} = {-}{0.05em},
    }
    \textbf{Name} & \textbf{Description} & \textbf{Data Type} \\
    polarization & Acquired polarizaions & string \\
    acquisition\_date & Image acquisition dates, formatted as YYYY-MM-DD & string \\
    x & UTM/UPS x (East) coordinates [m] & float64 \\
    y & UTM/UPS y (North) coordinates [m] & float64 \\
\end{longtblr}

\subsubsection{Main Metadata}
\label{sec:enh_product_metadata}
The following table lists the main metadata and per-pixel metadata associated with the \acs{enrb} and  \acs{eorb} products.

\begin{longtblr}[
    label = tab:meta_enh_products,
    entry = none,
    caption = { \acs{enrb} and  \acs{eorb} main metadata and per-pixel metadata},
    ]{
        width = \linewidth,
        colspec = {Q[l,25]Q[l,38]Q[l,10]Q[l,10]Q[c,12]},
        hlines,
        hline{1-Z} = {-}{0.05em},
    }
    \textbf{Name} & \textbf{Description} & \textbf{Dimensions} & \textbf{Data Type} & \textbf{Per-pixel Metadata}\\
    reference\_slc / reference\_cslc & List of links to input reference \acs{slc} or \acs{cslc} & - & list of string & \\
    reference\_dem & Link to reference DEM. Only for \acs{enrb} & - & string & \\
    reference\_geoid & Link to reference GEOID. Only for \acs{eorb} & - & string & \\
    etad\_file & ETAD file used for each image belonging to the \acs{cslc}/\acs{slc} stack. This is optional: it is included only if used; otherwise, it is null & - & string & \\ 
    tec & Link to external \acs{tec} data, if used. & - & string & \\ 
    era5 & Link to external \acs{era5} data, if used.  & - & string & \\
    reference\_inc\_angle & Optional attribute: list of link to reference DEM- or GEOID-based local incidence angle maps. It is present only if the \acs{nrb} is generated from \acs{cslc}s. & - & string & \ding{51} \\
    ground\_heading & Azimuth direction. Only for \acs{eorb} & x, y & float32 & \ding{51} \\
    reference\_anf & Optional attribute: list of links to reference area normalization factor or local contributing area; only for \acs{enrb}. It is present only if the \acs{enrb} is generated from \acs{cslc}s. & - & string & \\
    orig\_data\_id & Per pixel metadata reporting an index describing the bursts contributing to the output mosaicked pixel & x, y & int32 & \ding{51} \\
    mask & Mask containing, for each ground cell, the following information:
    
     - 0: no distortions
    
     - 1: layover
    
     - 2: shadow
    
     - 3: undefined.

     Only for \acs{enrb}
     & x, y & byte & \ding{51} \\
    nlooks & Number of equivalent looks for the gamma flattening & x, y & float32 & \ding{51} \\ 
    \_geometric\_accuracy & Measured geometric accuracy, including eastern and northern geometric accuracy bias and standard deviation (only for \acs{enrb}) & - & dictionary & \\ 
    east\_geom\_bias & Measured eastern geometric bias (only for \acs{enrb}) & - & float32 & \\
    east\_geom\_std & Measured eastern geometric standard deviation (only for \acs{enrb}) & - & float32 & \\
    north\_geom\_bias & Measured northern geometric bias (only for \acs{enrb}) & - & float32 & \\
    north\_geom\_std & Measured northern geometric standard deviation (only for \acs{enrb}) & - & float32 & \\
    geometric\_accuracy
    \_radial\_rmse& Measured radial RMSE of geometric accuracy (only for \acs{enrb}) & - & float32 & \\
    x\_sample & Pixel spacing in x direction & - & float32 & \\
    y\_sample & Pixel spacing in y direction & - & float32 & \\
    ghost\_rm\_flag & Flag indicating the application of ghosts removal & - & string & \\
    ghost\_rm
    
    \_algo & Link to the algorithm used for ghosts removal & - & string & \\
    side\_lobe
    \_reduction\_flag & Flag indicating the application of side lode reduction & - & string & \\
    side\_lobe
    
    \_reduction
    
    \_algo & Link to the algorithm used for side lobe reduction & - & string & \\
    
\end{longtblr}

The table below describes the dimensions/coordinates associated with the \acs{enrb} and \acs{eorb} main metadata.

\begin{longtblr}[
    label = tab:var_enh_metadata_coordinates,
    entry = none,
    caption = {\acs{enrb} and \acs{eorb} main metadata coordinates/dimensions},
    ]{
        width = \linewidth,
        colspec = {Q[l,20]Q[l,62]Q[l,18]},
        hlines,
        hline{1-Z} = {-}{0.05em},
    }
    \textbf{Name} & \textbf{Description} & \textbf{Data Type} \\
    x & UTM/UPS x (East) coordinates [m] & float64 \\
    y & UTM/UPS y (North) coordinates [m] & float64 \\
\end{longtblr}

\subsubsection{Reference System Tiling and Slicing}
\label{sec:enh_reference_system_tiling}

\autoref{tab:enh_ref_system_slicing} reports the reference system, the slicing, and grid for \acs{enrb} and \acs{eorb} products.

\begin{longtblr}[
    label=tab:enh_ref_system_slicing,
    entry = none,
    caption = {Reference system and Slicing for \acs{enrb} and \acs{eorb}},
    ]{
        width=\linewidth,
        colspec = {Q[l,15]Q[l,35]Q[l,35]Q[l,15]},
        %vline{2-5} = {-}{},
        %vline{5-5} = {-}{},
        hlines,
        hline{1-Z} = {-}{0.05em},
        rowhead=1
    }
    \textbf{Product} & \textbf{Reference System} & \textbf{Slicing} & \textbf{Spacing}\\
    \textbf{\acs{enrb}} (default) & Cartographic: UTM/UPS & US-\acs{mgrs} 100 km x 100 km & 10 m\\
    \textbf{\acs{enrb}} (optional) & Cartographic: UTM/UPS & US-\acs{mgrs} 100 km x 100 km & 30 m\\
    \textbf{\acs{eorb}} (default) & Cartographic: UTM/UPS & US-\acs{mgrs} 100 km x 100 km & 10 m\\
    \textbf{\acs{eorb}} (optional) & Cartographic: UTM/UPS & US-\acs{mgrs} 100 km x 100 km & 30 m\\	
\end{longtblr}

\subsection{Product Requirements}
No additional criticalities are foreseen beyond those listed for standard \acs{nrb} or \acs{orb} products.

\subsubsection{Accuracy}

% La risoluzione 

As regards Ghost Removal, two approaches have been evaluated (both the analysed methods base on the asymmetry of the Doppler Spectrum of \acs{sar} ghosts):
\begin{itemize}
    \item The “asymmetric mapping and selective filtering” methods described in \cite{Avolio_et_al2014} and \cite{DiMartino_et_al2014} mitigates azimuth ambiguities through a two-step process: asymmetric Wiener filters are applied to isolate ambiguity-free spectral regions, and binary maps identify affected areas. Only those ambiguous regions are selectively replaced with their filtered counterparts, preserving the rest of the image.
    \item The method in \cite{Choi2021} exploits the spectral distinction between azimuth ambiguities and genuine targets: while true targets exhibit broad spectral support, ghost signals are narrowly concentrated near the Doppler Nyquist frequency. By forming three sub-aperture images and projecting them onto a rotated coordinate system, the method concentrates target power into a single component while spreading ghost energy across all components, enabling effective ambiguity suppression through simple recombination.
\end{itemize}
This technique is often seen as an intermediate step for further analyses, such as detection (especially in maritime environments) or classification, but some metrics can help assessing the performances. The following table summarizes the accuracy requirements for the Ghost Removal technique for \acs{enrb} and \acs{eorb} products, considering coastal environments as test site.

\begin{longtblr}[
	label = tab:enrb_ghost_accuracy,
	entry = none,
	caption = {Accuracy Requirements for \acs{enrb} / \acs{eorb} products, regarding Ghost Removal (Azimuth and Range Ambiguities)},
]{
	width = \linewidth,
	colspec = {Q[l,30] Q[l,60] Q[c,20]},
	hlines,
	hline{1-Z} = {-}{0.05em},
	rowhead = 1,
}

\textbf{Metric} & \textbf{Description} & \textbf{Requirement} \\

\textbf{ASR / RSR (dB)} & Ambiguity-to-Signal Ratio in azimuth (ASR) or range (RSR), defined as the ratio of average power in ambiguity regions (typically displaced by integer multiples of PRF or sampling interval) to the mainlobe peak power. & $\leq -28$~dB \\

\textbf{GBR Reduction (dB)} & Difference in ghost-to-background ratio before and after suppression. GBR is computed as the ratio between the average power in ambiguity-affected regions and the average power in nearby background (clutter or noise) regions. & $\geq 5$~dB \\

\end{longtblr}

As regards Sidelobe Suppression techniques, Spatially Variant Apodization (SVA) technique (and its refinements) have been considered, due to their capacity to suppress sidelobes while preserving the resolution for strong scatterers, following \cite{Fischer2006} and \cite{Liu2018}. The following table summarizes the accuracy requirements for the Sidelobe Suppression technique for \acs{enrb} and \acs{eorb} products, considering a strong pointlike scatterer in a uniform region as test site.

\begin{longtblr}[
	label = tab:enrb_sidelobe_accuracy,
	entry = none,
	caption = {Accuracy Requirements for \acs{enrb} / \acs{eorb} products, regarding Sidelobe Suppression application},
]{
	width = \linewidth,
	colspec = {Q[l,30] Q[l,60] Q[c,20]},
	hlines,
	hline{1-Z} = {-}{0.05em},
	rowhead = 1,
}

\textbf{Metric} & \textbf{Description} & \textbf{Requirement} \\

\textbf{PSLR (dB)} & Peak sidelobe level relative to mainlobe peak & $\leq -25$~dB \\

\textbf{ISLR (dB)} & Total sidelobe energy relative to mainlobe energy & $\leq -10$~dB \\

\textbf{FWHM Deviation (\%)} & Full Width at Half Maximum variation from theoretical mainlobe width & $\leq 10\%$ \\

\textbf{SCNR (dB)} & Signal-to-clutter-plus-noise ratio in target region & $\geq 10$~dB \\

\textbf{Sidelobe Floor (dB)} & Average level of sidelobe floor excluding peak sidelobes & $\leq -30$~dB \\

\end{longtblr}

\subsubsection{Timeliness}

\acs{enrb} and \acs{eorb} products can be generated as soon as L1 data are available, for each \acs{slc} image independently or starting from \acs{cslc} products, with the addition of the Enhancing step.

The Enhancing step can be performed on \acs{slc} or \acs{cslc} products, and does not require any other input data. Additionally, for the timeliness of the identified algorithms, the following considerations are valid:
\begin{itemize}
    \item Ghost Removal: the increment of timeliness is negligible with respect to standard \acs{nrb} or \acs{orb} production.
    \item Sidelobe Suppression: it is of the same order as that required to generate the corresponding standard products. 
\end{itemize}
The composition of the Enhancing step is different for \acs{enrb} and \acs{eorb} products, but considerations about timeliness are the same. The only difference involves Ghost Removal, which is an operation useful only in the maritime environment.

In the framework of this project, the \acs{cslc} procedure includes geometric correctness, even in case of parameter compensation that does not introduce differential errors between reference and secondary images. Taking into account this characteristic, if the \acs{enrb} is generated from the \acs{cslc}, all the calibrations are already performed during the \acs{cslc} processing, and the reference geometry is computed only once for the reference orbit. Therefore, the timeliness does not change depending on the information needed to perform the Geo-to-Radar conversion or to generate the \acs{cslc} product.

The main factor impacting the timeline is instead the availability of external data or reference images, which are required to refine geolocation and mitigate atmospheric effects, particularly ionospheric effects. 

\autoref{fig:enrb_single} shows the flowchart for generating the \acs{enrb}/\acs{eorb} product using a single \acs{slc} image.

\begin{figure}[H]
	\centering
	\resizebox{\textwidth}{!}{
		\begin{tikzpicture}[node distance=0.5cm and 2cm]
			
			% Define the style for the nodes
			% \tikzstyle{node} = [circle, draw, fill=blue!20, text centered, minimum size=1.5cm]
			
			% Draw nodes
			\node[rectin] (init) {\acs{slc}};
                \node[block, right=of init] (enh) {Enhancing};
			\node[block, right=of enh] (g2r) {Geo-to-Radar};
			\node[block, right=of g2r] (rtc) {\acs{rtc}};
			\node[rectaux, below=of init] (dem) {\acs{dem}};
			\node[rectaux, below=of dem] (orbit) {Orbital data};
			\node[rectauxd, below=of orbit] (models) {External models};
			\node[rectind, below=of models] (ref_stack) {Reference SLC stack};
			\node[block, below=of rtc] (geocoding) {Geocoding};
                \node[rectout, below=of geocoding] (nrb) {\acs{enrb}/\acs{eorb}};
			
			% Draw edges with arrows
			\draw[->, thick] (init) -- (enh);
                \draw[->, thick] (enh) -- (g2r);
			\draw[->, thick] (g2r) -- (rtc);
			\draw[->, thick] (rtc) -- (geocoding);
			\draw[->, thick] (dem) -| (g2r);
			\draw[->, thick] (orbit) -| (g2r);
			\draw[->, thick] (models) -| (g2r);
			\draw[->, thick] (ref_stack) -| (g2r);
			\draw[->, thick] (geocoding) -- (nrb);
		\end{tikzpicture}
	}
	\caption{Sample workflow for generation of \acs{enrb} considering each \acs{slc} image separately.}
	\label{fig:enrb_single}
\end{figure}

\autoref{fig:enrb_from_cslc} shows the flowchart for generating the \acs{enrb}/\acs{eorb} product using \acs{cslc} products.

\begin{figure}[H]
	\centering
	\resizebox{\textwidth}{!}{
		\begin{tikzpicture}[node distance=0.5cm and 2cm]
			
			% Define the style for the nodes
			% \tikzstyle{node} = [circle, draw, fill=blue!20, text centered, minimum size=1.5cm]
			
			% Draw nodes
			\node[rectin] (init) {\acs{slc}};
            \node[block, right=of init] (coreg) {Co-registration};
            \node[rectout, right=of coreg] (cslc) {\acs{cslc}};
            \node[block, right=of cslc] (enh) {Enhancing};
            \node[block, below=of enh] (rtc) {\acs{rtc}};    
            \node[block, below=of rtc] (geocoding) {Geocoding};
    
            \node[rectaux, below=of init] (dem) {\acs{dem}};
            \node[rectaux, below=of dem] (orbit) {Orbital data};
            \node[rectauxd, below=of orbit] (models) {External models};
            \node[rectind, below=of models] (ref_stack) {Reference SLC stack};
            \node[rectout, below=of geocoding] (enrb) {\acs{enrb}/\acs{eorb}};
    
            % Draw edges with arrows
            \draw[->, thick] (init) -- (coreg);
            \draw[->, thick] (coreg) -- (cslc);
            \draw[->, thick] (cslc) -- (enh);
            \draw[->, thick] (enh) -- (rtc);
            \draw[->, thick] (rtc) -- (geocoding);
    
            % Draw edges with arrows
            \draw[->, thick] (dem) -| (coreg);
            \draw[->, thick] (orbit) -| (coreg);
            \draw[->, thick] (models) -| (coreg);
            \draw[->, thick] (ref_stack) -| (coreg);
            \draw[->, thick] (geocoding) -- (enrb);
        \end{tikzpicture}
	}
	\caption{Sample workflow for generation of \acs{enrb} starting from \acs{cslc} products.}
	\label{fig:enrb_from_cslc}
\end{figure}

\subsection{Usability, Integrability and Interoperability}
\label{sec:enrb_usability_integrability}

The \acs{enrb} and \acs{eorb} products are improved versions of the \acs{nrb} and \acs{orb} products respectively. For this reason, Sections \autoref{sec:nrb_usability_integrability} for \acs{enrb} and \autoref{sec:orb_usability_integrability} for \acs{eorb} describe already the applications for these products, the usability properties as well as the integrability into different systems and workflows.

Applying techniques, like Ghost Removal and Sidelobe Suppression, significantly improves the quality and interpretability of radar imagery, particularly in high-resolution or complex environments, leading to cleaner and more reliable radar products, that enable better feature extraction, classification, and change detection in SAR applications.

In some of the analysed \acs{nrb} products the speckle filter is not applied. For specific analyses like land cover classification or vegetation monitoring, speckle noise mitigation may be required. This is an additional preprocessing step that one may need to consider.

Finally, radar backscatter can be influenced by surface roughness, moisture content, and other environmental factors. These must be accounted for when interpreting the data, especially in applications where surface characteristics (e.g., soil moisture, vegetation growth) vary over time.